% EigenKache first-paper scaffold. Replace every TODO with measured evidence.
\documentclass[11pt]{article}
\usepackage[margin=1in]{geometry}
\usepackage{amsmath,booktabs,graphicx,hyperref}
\title{EigenKache: Attention-Conditioned Landmark Compression for Memory-Bounded Long-Context Decoding}
\author{Aryan Putta}
\date{\today}
\begin{document}
\maketitle

\begin{abstract}
\textbf{TODO:} State the memory problem, method, evaluation, main measured result, and limitation in 150--200 words. Do not include an unsupported claim.
\end{abstract}

\section{Introduction}
\textbf{TODO:} Explain why long-context decoding creates KV-cache pressure. End with the research question and three contributions.

\section{Background and Related Work}
\textbf{TODO:} Explain attention, KV caching, FlashAttention's IO perspective, paging, eviction, and compression. Compare methods by what they preserve and what they cost.

\section{Method}
\textbf{TODO:} Define sink tokens, hot tail, cold middle, attention-conditioned segments, weighted landmarks, and the token budget. Include pseudocode and a diagram.

\section{Experimental Setup}
\textbf{TODO:} Document model, traces, datasets, hardware, software versions, seeds, budgets, baselines, and metrics so another person can reproduce the tables.

\section{Results}
\textbf{TODO:} Report quality/error versus budget, memory savings, policy overhead, and GPU runtime only when measured on GPU hardware.

\section{Ablations and Limitations}
\textbf{TODO:} Vary sink/tail sizes, segment count, head behavior, sequence length, and workload. State what the current evidence cannot establish.

\section{Conclusion}
\textbf{TODO:} Answer the research question narrowly and identify the next experiment.

\bibliographystyle{plain}
% TODO: add primary papers and official dataset/model citations.
\end{document}
